\documentclass[12pt]{exam}
\usepackage[utf8]{inputenc}

\usepackage{amsmath,amstext,amsthm,amssymb,amsxtra}
\usepackage[top=1.5in, bottom=1.5in, left=1.25in, right=1.25in]	{geometry}
%\usepackage[normalem]{ulem}
\usepackage{txfonts} % pxfonts txfonts 
\usepackage[T1]{fontenc}
\usepackage{lmodern}
\renewcommand*\familydefault{\sfdefault}
 \usepackage{euler}   % better than the option below
\usepackage{pdfsync}
\usepackage{multicol}
\newcommand{\ci}[1]{_{ {}_{\scriptstyle #1}}}
\usepackage{graphicx}
\graphicspath{ {images/} }


\newcommand{\norm}[1]{\ensuremath{\left\|#1\right\|}}
\newcommand{\abs}[1]{\ensuremath{\left\vert#1\right\vert}}
\newcommand{\ip}[2]{\ensuremath{\left\langle#1,#2\right\rangle}}
\newcommand{\p}{\ensuremath{\partial}}
\newcommand{\pr}{\mathcal{P}}

\newcommand{\pbar}{\ensuremath{\bar{\partial}}}
\newcommand{\db}{\overline\partial}
\newcommand{\D}{\mathbb{D}}
\newcommand{\B}{\mathbb{B}}
\newcommand{\Sp}{\mathbb{S}}
\newcommand{\T}{\mathbb{T}}
\newcommand{\R}{\mathbb{R}}
\newcommand{\Z}{\mathbb{Z}}
\newcommand{\C}{\mathbb{C}}
\newcommand{\N}{\mathbb{N}}
\newcommand{\Q}{\mathbb{Q}}
\newcommand{\mQ}{\mathcal{Q}}
\newcommand{\mS}{\mathcal{S}}
\newcommand{\scrH}{\mathcal{H}}
\newcommand{\scrL}{\mathcal{L}}
\newcommand{\td}{\widetilde\Delta}
\newcommand{\pw}{\text{PW}}
\newcommand{\esup}{\text{ess.sup}}
\newcommand{\Tn}{\mathcal{T}_n}
\newcommand{\Bn}{\mathbb{B}_n}
\newcommand{\rt}{\mathcal{O}}
\newcommand{\avg}[1]{\langle #1 \rangle}
\newcommand{\one}{\mathbbm{1}}
\newcommand{\eps}{\varepsilon}
\newcommand{\grad}{\nabla}

\newcommand{\La}{\langle }
\newcommand{\Ra}{\rangle }
\newcommand{\rk}{\operatorname{rk}}
\newcommand{\card}{\operatorname{card}}
\newcommand{\ran}{\operatorname{Ran}}
\newcommand{\osc}{\operatorname{OSC}}
\newcommand{\im}{\operatorname{Im}}
\newcommand{\re}{\operatorname{Re}}
\newcommand{\tr}{\operatorname{tr}}
\newcommand{\vf}{\varphi}
\newcommand{\f}[2]{\ensuremath{\frac{#1}{#2}}}

\newcommand{\kzp}{k_z^{(p,\alpha)}}
\newcommand{\klp}{k_{\lambda_i}^{(p,\alpha)}}
\newcommand{\TTp}{\mathcal{T}_p}
\newcommand{\m}[1]{\mathcal{#1}}
\newcommand{\md}{\mathcal{D}}
\newcommand{\qan}{\abs{Q}^{\alpha/n}}
\newcommand{\sbump}[2]{[[ #1,#2 ]]}
\newcommand{\mbump}[2]{\lceil #1,#2 \rceil}
\newcommand{\cbump}[2]{\lfloor #1,#2 \rfloor}

\newcommand{\class}{Math 629 Spring 26}
\newcommand{\term}{Fall 24}
\newcommand{\examnum}{Homework \hwnum}
\newcommand{\examdate}{\duedate}
\newcommand{\timelimit}{75 Minutes}
\newcommand{\vc}[3]{\langle #1,#2,#3\rangle}
\newcommand*{\vv}[1]{\vec{\mkern0mu#1}}
\newcommand{\bv}[1]{\boldsymbol{#1}}
\newcommand{\hide}[1]{}
\newcommand{\uvec}[1]{\boldsymbol{\hat{\textbf{#1}}}}
\newcommand{\vex}[1]{\boldsymbol{{\textbf{#1}}}}

\pagestyle{head}
\firstpageheader{}{}{}
\runningheader{\class}{\examnum\ - Page \thepage\ of \numpages}{\examdate}
\runningheadrule

\makeatletter
\renewcommand*\env@matrix[1][*\c@MaxMatrixCols c]{%
  \hskip -\arraycolsep
  \let\@ifnextchar\new@ifnextchar
  \array{#1}}
\makeatother

%\printanswers
\begin{document}

\noindent
\begin{tabular*}{\textwidth}{l @{\extracolsep{\fill}} r @{\extracolsep{6pt}} l}
\textbf{\class} & \textbf{Name: Mengxiang Jiang}\\
\end{tabular*}\\
\rule[2ex]{\textwidth}{2pt}
%

 

\newcommand{\duedate}{02/23/2026 at 11:59pm}
\newcommand\hwnum{6}

Turn this assignment in on gradescope. All work must be neat and legible. Turn in 
a ``final draft''. There should be nothing scratched out and your work should flow
generally from left to right and top to bottom.

(Yes - there are only two problems!)

\begin{questions}
\question 
Find two non--trivial solutions to $x^2 - 60y^2 = 1$. (Use the 
Chakravala algorithm; you don't need to show the arithmetic, but you
do need to show the steps in the algorithm. That is, the 
$(x,y,k)$ in each step).
\\\\
Since at each step of the algorithm, we want to minimize $|m^2-N|$,
we start with $m=8$ since $8^2 = 64$ is closest to 60,
 and get the first triple $(8,1,4)$. For our first improvement,
 we need to find the next $m$ such that $\frac{x+ym}{k}$ is an integer, 
 while still minimizing $|m^2-60|$. We see that $m \equiv 0 \mod 4$, and
 again, we have $m=8$ is the closest to $\sqrt{60}$. Our next $x$ is 
 \[
  x_{new} = \frac{xm + 60y}{|k|} = \frac{8*8 + 60*1}{4} = 31.
 \]
 Our next $y$ is
 \[
  y_{new} = \frac{x + ym}{|k|} = \frac{8 + 1*8}{4} = 4.
\]
Our next $k$ is
\[
  k_{new} = \frac{m^2 - 60}{k} = \frac{64 - 60}{4} = 1.
\]
Thus, we have found one solution, namely $(31, 4, 1)$.
Now, we can pick a different starting $m$ for our first step. 
We can pick $m=7$ since $7^2 = 49$ is the second closest to 60, 
and get the first triple $(7,1,-12)$. I will omit the arithmetic and just list
the steps:
\begin{center}
  \begin{tabular}{|c|c|c|c|}
    Best $m$ & $x_new$ & $y_new$ & $k_new$\\
    \hline
    5 & 7 & 1 & 2 \\
    7 & 54 & 7 & -6 \\
    6 & 124 & 16 & 4 \\
    8 & 488 & 63 & 1\\
  \end{tabular}
\end{center}
\newpage 
\question 
Using the algorithm with starting values $(10,1,8)$ find a solution to
$x^2 - 92y^2 = 1$. Once you get to the point where $k=\pm4$, use some 
reduction to get to a solution (note: I showed how to do this when 
$k = \pm 2$; do something similar here). Show your work as I mentioned 
above. 
\begin{solutionorbox}[\stretch{1}]
 
\end{solutionorbox}



\end{questions}
\end{document}
